\subsection*{领取发布工作以后再发送}

EV1 已经提交，P1 为 \code{paid}，平台余额为 7000。发布器 W1 扫描 \code{phase=ready} 且 \code{next\_run\_at} 已到的事件，也扫描领取已到期的事件。扫描提供候选编号，领取事务锁住 EV1，重新核对当前状态，再用数据库锁后采样的时刻判断资格。这里继续采用 PostgreSQL 18 的 Read Committed 与短事务行锁，并沿用已确认本地提交在所讨论故障后继续保留的前提~\cite{pg}。

允许领取时，事务递增 \code{publish\_epoch}，保存本次进程运行身份与租期。资格判断、锁等待和真实时间的关系沿第八章规则：等待以后再采样，已经过期的旧领取重新取得新代次。事务先提交领取与尝试记录，发布器再按提交结果执行远程发送。

这次领取还保存发布尝试 PA1。\code{attempt\_id} 标识一次向代理发送的尝试，由发布器分配并在发送以前提交。它引用 EV1，记录本轮运行身份和代次。再次调用发布时，即使语义报文相同，也另建 PA2。支付调用凭 \code{pay-C1} 归并到原付款，消息副本凭 EV1 归并到原事件，尝试记录逐次保存各次通信的经过。

代理连接中的 channel 是一条具有独立协议状态的逻辑通道。RabbitMQ 的发布确认在每个开启确认模式的 channel 内计数，从 1 开始。序号与该 channel 配合，才指向确定的发布~\cite{rabbitmq-confirms}。本地给每次 channel 建立新的 \code{channel\_incarnation}，防止重连后序号 1 与旧通道的序号 1 混淆。发布器在同一通道的串行发送位置取得下一次实际发布序号，并保存对应尝试，提交确认以后才执行这次发送。登记与发送之间，该通道保持这个发送位置。登记提交未知时先查明，进程恢复则放弃旧通道，另建通道与尝试。尚未真正发送时，代理序号保持原值，适配器为已登记发布保留这一位置。

\begin{BookCode}
CREATE TABLE publish_attempts (
  attempt_id text PRIMARY KEY,
  event_id text NOT NULL REFERENCES outbox_events(event_id),
  owner_incarnation text NOT NULL,
  publish_epoch bigint NOT NULL,
  channel_incarnation text NOT NULL,
  publish_seq bigint NOT NULL,
  outcome text NOT NULL
    CHECK (outcome IN ('pending','accepted','returned','nacked','unknown')),
  receipt jsonb,
  UNIQUE (channel_incarnation,publish_seq)
);
\end{BookCode}

\code{outcome} 保存这次尝试已取得的通信结论，\code{receipt} 保存可恢复的确认或退回内容。发送以前为 \code{pending}，连接失效而结果未明时可以记为 \code{unknown}，尝试身份继续保留。没有收到确认只说明本地缺少证据，代理中可能已有副本。

若领取事务提交后、真正发送前进程退出，EV1 与 PA1 都在，代理中还没有消息。租期结束以后，W2 领取新代次，保存 PA2，再发送原 EV1。PA1 的记录解释了前一次准备过程，EV1 则继续留在待完成范围。

\subsection*{代理受理与本地保存之间的两次提交}

现在沿一次完整发布安排故障位置。W1 读取 EV1 的原报文与目标配置，把内容写入连接。写入本地 socket 完成，只说明本地网络栈接收了这些字节。代理可能尚未收到、尚未处理，或者稍后才收到~\cite{rabbitmq-confirms}。

代理在自己的受理条件满足以后发回确认。W1 收到并核对该确认，再用一个本地事务保存 PA1 的证据，并把 EV1 改成 \code{accepted}。这个状态在本章的含义是：指定目标在所采用的持久性前提下受理了至少一个副本。履约系统的业务完成情况由它自己的处理记录表达。

\CNineFigure{ch09-original04}{发布从发送前到本地受理记录提交后的故障位置}{fig:ch09-publish-failures}

在代理受理以前退出，W2 重发可能形成第一个副本。在代理已经受理以后、确认尚未到达时断开，W2 看到缺少证据，重发会形成第二个副本。确认已经到达 W1 内存、保存本地记录以前退出，也有同样结果。即使本地 \code{UPDATE} 已经执行，只要保存确认的事务尚未提交，退出后依然回到原待完成状态。

最后一处在本地受理记录提交以后。W1 丢失该事务的返回值，重读 EV1 可以看到 \code{accepted} 和对应证据，便结束本轮。若只用内存布尔值记已发布，进程重启就要重新面对缺少证据的局面。

发布器在取得受理证据以后停止重试，就要容纳中间几种历史中的重复。设 H0 的 PA1 没到代理，H1 的 PA1 已被代理接收但确认丢失。本地 EV1、PA1 和所见连接错误可以完全相同。H0 要靠重发继续，相同恢复决定在 H1 中可能再建一个副本。稳定 EV1 使两份副本指向同一事实~\cite{outbox,rabbitmq-reliability}。

\CNineFigure{ch09-unknown-histories}{同一份本地未知状态对应两条代理历史，恢复重发分别形成第一或第二个副本}{fig:ch09-unknown-histories}

把 EV1 先标成 \code{accepted} 再发，可以减少这类重发，却会产生标记已提交、发送未发生的漏通知。Outbox 保存待办，确认提供停止重试所需的证据。两次提交之间出现未知结果时，恢复者继续发布同一事件，消费者再按稳定身份识别其副本。

\subsection*{RabbitMQ 的确认究竟确认了什么}

以 RabbitMQ 4.3 对应本章的代理受理。发布者开启 publisher confirms，即由代理向发布者回报受理结论的确认模式。代理用 \code{basic.ack} 或 \code{basic.nack} 回报。前者是发布确认，后者表示代理未能成功处理该次发布并拒绝承担该消息的责任。发布 nack 中的 \code{requeue} 字段应忽略，消费端重新入队则有另一套处理规则~\cite{rabbitmq-confirms}。

路由信息也参与判断。消息没有匹配任何队列时，RabbitMQ 也可以发回发布确认。发布设置 \code{mandatory=true}，完全无法路由的消息先经 \code{basic.return} 退回，随后才收到对应 ack 或 nack~\cite{rabbitmq-confirms}。本地把退回与确认合在一起裁定。某次尝试先被退回，后来收到 ack，保存的结论就是该次发布已完成协议处理，但路由失败，事件回到重试或核查路径。

进一步看一个反例。原计划发送到履约队列，配置错误却只绑定到审计队列。消息进入审计队列，代理可以 ack，也不会 return，因为确实存在匹配队列。履约系统却没有副本。\code{mandatory} 提供的是至少存在一个路由目标的检查，指定履约目标还依赖被核对并持续受控的拓扑~\cite{rabbitmq-reliability}。

本章直接向已声明的 \code{fulfillment.paid} 队列发送：使用默认交换机，以该队列名作为路由键，并核对同一虚拟主机中的实际队列类型、持久性与存续策略。虚拟主机划分代理中的命名与权限范围，默认交换机则按队列名自动建立路由绑定~\cite{rabbitmq-exchanges}。发布身份按权限限制在这个目的范围，队列名的重建、删除和配置变化由受控管理流程处理。改用自定义交换机和绑定时，管理流程也要保证当前路由集合包含原定履约队列。

受控管理覆盖核对以后到受理期间。管理流程保持原目标存续。确需迁移时，先停止相关新发布，处理旧通道与未知尝试，再为原事件保存指向同一业务目的的迁移依据。若发送前查过绑定以后就允许任意删除队列，队列可能恰好在检查与发送之间消失。迁移协议要把这一段变化纳入切换决定。

持久消息路由到持久队列时，RabbitMQ 的确认条件涉及队列持久保存。quorum 队列是以多个成员复制队列历史的队列类型，在多数副本接受并向 leader 确认以后发出相应发布确认~\cite{rabbitmq-confirms}。本章采用持久消息与持久的 quorum 履约队列，成员至少为三个。已确认消息的安全依赖多数成员没有永久失去，继续运行则要求足够多数恢复可用~\cite{rabbitmq-quorum-queues}。目标队列的副本条件与本地 PostgreSQL 的持久条件分别成立。

发布确认与消费者确认各管一段通信。\footnote{RabbitMQ 4.3 文档称两类确认“entirely orthogonal and unaware of each other”。发布确认依据代理接收，消费确认依据消费端选定的完成条件。第十章把该完成条件接到业务事务~\cite{rabbitmq-confirms}。}发布 ack 已经到达，消费者可能尚未运行。消费者也可能先收到消息，发布者的确认还在途中。EV1 的 \code{accepted} 表达发布边界已经取得所需证据。

\subsection*{怎样把确认送回正确的尝试}

PA1 与 PA2 都发布 EV1 时，仅按 \code{event\_id} 查找最新尝试会误配：PA1 的退回可能被算到 PA2，PA2 的确认又可能覆盖 PA1 的失败。每次发送把唯一 \code{attempt\_id} 放入原消息属性，例如自定义属性 \code{x-publish-attempt}，语义报文继续携带原 EV1。

\code{basic.return} 带回原消息内容与属性~\cite{amqp091}。适配器读取其中的 PA1，核对原通道及保存的报文，再把退回归到 PA1。发布确认携带通道内序号，适配器用 \code{(channel\_incarnation,publish\_seq)} 找到同一 PA1。这两条映射在发送前建立。通道内发送、退回和确认按同一适配器的有序事件循环处理，先记录退回，再裁定随后确认。

\CNineFigure{ch09-return-confirm}{退回按消息属性找尝试，确认按通道与序号找尝试，再合并裁定}{fig:ch09-return-confirm}

确认还可能乱序，或者以 \code{multiple=true} 一次覆盖到某序号为止的多项发布~\cite{rabbitmq-confirms}。例如本通道有序号 1、2、3，先确认 3，再收到覆盖至 2 的批量确认。适配器分别关闭 3，再关闭尚未完成的 1、2。若 2 先有退回，它保持路由失败。另一通道的序号 1 属于另一条记录。消费投递的 delivery tag 也只有 channel 内作用域，它用来确认特定投递，业务事件则用 EV1 识别。

通道失效时，适配器逐项保存尚未确定的发布。PA1 若缺少完整确认，就记为 \code{unknown}，保留原通道、序号和已收到的观察。恢复者建立新通道，取得新的运行身份和序号映射，再为重发建立 PA2、PA3。已判定的 PA 再次收到相同观察，读取原证据即可。若同一 PA 的两份观察互相矛盾，适配器保留两份内容，登记协议或关联核查，由核查者沿原映射查明差异。

\subsection*{接管以后到来的确认怎样处理}

W1 持有发布第 7 代，PA1 已送出。租期结束，W2 接管第 8 代并准备 PA2。此时 PA1 的成功受理确认才到 W1，或者负责连接的适配器才恢复对它的处理。若仅凭 W1 已失去领取资格丢弃这项可信证据，恢复者就可能多发一个副本。

本章把发布调度与受理证据合流分开。W1 请求推迟下次运行、续期或把事件转回 \code{ready}，必须锁 EV1，匹配当前运行身份、代次、\code{leased} 和锁后未到期条件。PA1 的迟到失败只能更新它自己的尝试记录，EV1 的调度继续由 W2 当前领取决定。

PA1 的正面证据经独立入口核验原尝试、通道序号、原消息属性、无退回结论、目标及拓扑版本。入口锁 EV1，再锁 PA1，按同一事务保存证据并合流。只要证据满足原定受理条件，即使 PA1 的租期已过，也可以令 EV1 为 \code{accepted}，清除当前持有者并推进一次发布代次。W2 随后重新核对时会看到已经完成。

\CNineFigure{ch09-original05}{迟到的成功确认可以完成同一事件，迟到失败受当前调度资格限制}{fig:ch09-late-evidence}

W2 若已经发出 PA2，在途副本可以继续到达并被代理受理，两次尝试都关联 EV1。EV1 若已处于人工核查，迟到成功确认就与核查单共同保存解决依据。相应核查项在同一事务中结束，让核查者读取到已经解决的结论。

\begin{BookCode}
接收受信适配器的尝试观察；按 EV → PA 顺序加锁
核对 PA 引用的 EV、通道运行身份、序号、属性和原报文
若 PA 已有同内容终态，恢复原证据
若同 PA 的观察矛盾，保留冲突并登记核查
若为有效 ack，且本 PA 无 return、目标与持久条件满足：
    保存 PA.accepted 及完整受理证据
    EV 未 accepted 时：
        令 EV.accepted，保存该证据，清 owner/lease，发布代次加一
        存在核查单时共同结束它并记录解决依据
若为 return、nack 或结果未知：
    保存本 PA 的结论或未知观察
    仅当 EV 由本 PA 对应的运行身份与代次有效持有时，
        才能重新安排 EV.ready 或转入有归属的 EV.check
每项预期写入都要求恰好一行；异常或零行整笔回滚后重读
共同提交 PA 证据、EV 迁移及核查决定
\end{BookCode}

受理记录提交确认丢失时，重新读取 EV 与 PA。完整证据已经提交就恢复，尚未提交则继续原映射或等待接管。\code{accepted} 一经建立，后续失败观察只增加对应 PA 的记录。确认只在内存中存在、尚未持久保存时，恢复者根据未知结果继续，可能产生重复副本。
